% ============================================================
% RESUMEN
% ============================================================

\begin{abstract}

Este informe documenta cuatro ejercicios de procesamiento digital de señales desarrollados en el proyecto \texttt{WORK\_FOUR} (paquete \texttt{signal\_tools}). En el primero se cargó y segmentó la señal registrada en \texttt{Muestra01.csv} ($f_s = 5000$ Hz, 1000 muestras, cuatro segmentos de $0.05$ s) y se calculó la frecuencia dominante de cada segmento mediante FFT. En el segundo se generó una señal seno de $100$ Hz, se digitalizó con tres frecuencias de muestreo ($70$, $500$ y $1000$ Hz), se cuantizó con un algoritmo DAC de 3 bits, y se ilustró el fenómeno de aliasing y su prevención. En el tercero se descargaron grabaciones reales de un violín, un tambor y un gato, se calculó su FFT y se evaluó si el análisis de frecuencias permite diferenciarlos. En el cuarto se diseñó firmware mínimo para un ESP32-WROOM que compara tres mecanismos de muestreo periódico (bucle bloqueante de un solo hilo, tarea en el núcleo 0 con el núcleo 1 libre, y muestreo por interrupción de temporizador de hardware) frente a una carga de cómputo variable (suma de $1$ a $10\,000\,000$, $1\,000\,000$ y $100\,000\,000$); al no disponerse de hardware físico, esta sección se apoya en un modelo de temporización documentado y en una simulación etiquetada explícitamente como tal, en lugar de datos de serie reales. Todo el código, las pruebas automatizadas y los resultados de este informe están disponibles en el repositorio público \url{https://github.com/DanielAraqueStudios/Signals-frecuency}.

\end{abstract}

% ============================================================
% PALABRAS CLAVE
% ============================================================

\begin{IEEEkeywords}

Muestreo, aliasing, cuantización, Transformada Rápida de Fourier, FFT, ESP32, FreeRTOS, interrupciones de hardware, procesamiento digital de señales.

\end{IEEEkeywords}
